% ===== 原文 pp.107–111（PDF p.117–121） =====
\chapter{二维 \texorpdfstring{$\sigma$}{σ}-模型的缺陷}

% [原文 p.107]
众所周知，具有不同数量左移与右移超对称性的二维 $\sigma$-模型（例如 (0,1) 与 (0,2) 超对称 $\sigma$-模型）无法协调地定义在带非平凡边界的世界面上。这本质上是因为边界条件必须把左移与右移费米子的某个线性组合置零。

然而这并不排除把这些模型定义在带有缺陷（defect）的世界面上的可能性，所谓缺陷即一维子流形 $D$，两个不同 CFT 的场沿 $D$ 协调地粘合在一起。（边界条件是其中一个 CFT 为平凡的特例。）我们的目标是研究 (0,1) 与 (0,2) $\sigma$-模型中的超对称缺陷。

二维 CFT 的缺陷近来引起了很大兴趣，尤其是在有理共形场论的语境中，以及作为优雅地实现对偶性的手段 [55]、[53]、[51]、[52]。本章中我们将分析保持超共形对称性的判据，而不考虑这些理论可能具有的增强手征对称性；此外，我们将聚焦于缺陷的靶空间几何，只简要评论它们在实现对偶性中的作用。

考虑一个二维世界面，它被缺陷线 $D$ 分成两个区域 $\Sigma$ 与 $\hat\Sigma$。在区域 $\Sigma$ 上定义映射 $\Phi:\Sigma\to X$ 的 $\sigma$-模型，在区域 $\hat\Sigma$ 上定义映射 $\hat\Phi:\hat\Sigma\to\hat X$ 的 $\sigma$-模型，其中 $X$ 与 $\hat X$ 是两个紧致黎曼靶空间。这些映射到 $D$ 上的限制给出乘积映射 $\Phi\times\hat\Phi|_D:D\to X\times\hat X$。正如文献 [50] 所观察到的那样，粘合条件将要求该乘积映射取入靶空间乘积的某个子流形 $Y\subseteq X\times\hat X$ 之中。
% [Tutor Note: p.107 原文 “require the product map to takes values” 有笔误，应为 “to take values”（require sb. to do 的不定式不加 -s），译文按正确语法处理；$Y$ 的具体角色即 §9.2 的 Dirichlet 条件 $y\in Y$]

还可以在作用量中加入一项，让体场与活在 $Y$ 的世界体上的带联络线丛耦合（给它配上一个闭 2-形式 $F$），这完全类似于 $\sigma$-模型边界条件情形下膜上的 Chan-Paton 丛。借用文献 [50] 的术语，我们把对 $(Y,F)$ 称为双膜（bibrane）。

为保证缺陷理论保持体理论的某组指定对称性，我们
% [原文 p.108]
对嵌入 $Y$ 的几何以及 2-形式 $F$ 的选取施加某些要求。让我们通过列出相应的 Noether 荷来描述我们希望保持的那组对称性。(0,1) 超对称代数在 1+1 维中为
\[
\{Q_+,Q_+\}=H+P,
\]
其中 $Q_+$ 是唯一的、实的右移超荷，$H$ 是世界面能量，$P$ 是世界面动量。我们希望求出 $(Y,F)$，使得 $H$、$P$ 与 $Q_+$ 仍然是缺陷理论的守恒荷；由于 $H$ 与 $P$ 将分别守恒，我们要研究的将是文献中称为拓扑缺陷的那些例子 [52]。
% [Tutor Note: 这里的逻辑是：$\{Q_+,Q_+\}=H+P$ 只保证 $H+P$ 这一个组合守恒，而 p.110 的拓扑缺陷条件要求 $T^1{}_0$ 与 $T^1{}_1$ 分别连续粘合，从而 $H$ 与 $P$ 分别守恒——故“拓扑缺陷”严格强于“仅保持 (0,1) 超对称”；这一点在 §9.4 检验超流粘合时会被再次用到]

已经在玻色 $\sigma$-模型（即没有费米子场）的情形中，“哪些双膜 $(Y,F)$ 定义拓扑缺陷”就是一个有趣的问题，而这正是我们首先要关注的。结果发现，自然的选择是在（伪黎曼）乘积流形 $M=X\times\hat X$ 上取中性号差度规 $G=g\oplus-\hat g$，其中 $g$ 与 $\hat g$ 分别是 $X$ 与 $\hat X$ 上的正定度规。

那些给出拓扑缺陷的 $(Y,F)$ 与辛流形的 A-膜具有结构上的相似性（把环境靶空间上的辛形式换成上面定义的中性号差度规）；例如，正如 A-膜是关于靶空间辛形式的余迷向子流形一样，$Y$ 将被要求关于 $G$ 是“余迷向的”，其含义我们将在 §9.3 中解释。此外，我们将描述两类特殊的拓扑缺陷（“等距映射的图像”与“半 para-Kähler”缺陷），它们是拉格朗日 A-膜与填满空间的 A-膜的类比物。

另一种自然的做法是采用 Hitchin 广义几何 [58] 的语言，用它我们得到如下简单的刻画：$(Y,F)$ 定义玻色 $\sigma$-模型的拓扑缺陷，当且仅当 $Y$ 的“$F$-旋转广义切丛”被 $X\times\hat X$ 上的“广义度规”稳定化。我们将在 §9.3 中更深入地解释这些术语的含义。

我们把满足上述稳定化条件的 $(Y,F)$ 称为拓扑双膜；它们对任意中性号差流形 $M$ 都可定义，并且也许就其自身而言具有数学意义。然而我们应当指出，与当前物理讨论相关的中性号差流形属于非常受限的类型：即可以写成整体乘积 $M=X\times\hat X$、且度规限制到 $TX$（相应地 $T\hat X$）上为正（相应地负）定的流形。

在 §9.1 中我们一般地讨论缺陷粘合条件，并说明粘合条件保持体理论的一个对称性是什么意思。在 §9.2 中我们写下对应于 $(Y,F)$ 之选取的场粘合条件，并分析拓扑缺陷
% [原文 p.109]
要求。在 §9.3 中我们解释与 A-膜的类比，并用广义几何的语言重新表述拓扑双膜条件。在 §9.4 中我们分析 (0,1) 超对称 $\sigma$-模型的缺陷，在玻色粘合条件之外补充一个费米子粘合条件，并研究 $TY$ 的某个中间维子丛的几何。在 §9.5 中我们处理 (0,2) $\sigma$-模型的拓扑缺陷，将其描述为关于某个辛形式与复结构同时成为 A-膜与 B-膜的那些 $(Y,F)$。最后，在 §9.6 中我们简要评论我们的拓扑缺陷中实现 $\Sigma$ 与 $\hat\Sigma$ 上 $\sigma$-模型理论之间对偶性的一个子集。

本章正文中讨论的部分命题的证明将在附录 D 中给出。

\section{缺陷的粘合条件}

在讨论具体理论中的缺陷之前，我们先一般地讨论缺陷粘合条件，并说明缺陷保持体理论的一个对称性是什么意思。

作用量的变分一般由三类项组成：
\[
\delta S=(\delta S)_\Sigma+(\delta S)_{\hat\Sigma}+(\delta S)_D
\]
其中 $(\delta S)_\Sigma$ 与 $(\delta S)_{\hat\Sigma}$ 分别是区域 $\Sigma$ 与 $\hat\Sigma$ 上的积分，对满足体运动方程的场构形它们为零；而第三项 $(\delta S)_D$ 是 $D$ 上的积分，一般说来，除非我们沿 $D$ 对场施加粘合条件，它不会为零。粘合条件约束场沿 $D$ 的取值，从而也约束允许变分的集合，这里“允许变分”指把粘合条件的一个解映到另一个解的变分。因此，一个好的粘合条件的第一个要求是：它对所有允许变分都恒等地把 $(\delta S)_D$ 置零，并且是“极小地”做到这一点（即不过度约束沿 $D$ 的数据）。

此外，我们可能希望缺陷理论保持体的某个对称性。设 $(\sigma^0,\sigma^1)$ 是世界面坐标，在其中 $D$ 局部地由 $\sigma^1=0$ 的点集描述。在这些坐标下我们说一个粘合条件经典地保持体的一个对称性，当且仅当相应 Noether 流的 1-分量（即 $\sigma^1$ 分量）连续地穿过 $D$ 粘合，从而保证复合理论中存在守恒的 Noether 荷。（量子理论可能出现反常，但我们的讨论限于经典问题。）这一点满足之后，对称性变分将自动属于允许变分。
% [Tutor Note: p.109 的“1-分量”需要说清楚：世界面坐标取 $(\sigma^0,\sigma^1)$ 且 $D$ 为 $\sigma^1=0$，故 Noether 流 $J^\mu$ 的“1-分量”指沿 $\partial_1$ 的分量 $J^1$；要求 $J^1$ 在 $D$ 两侧连续粘合，等价于 $\int_{D} J^1\,\dd\sigma^0$ 的跃变为零，从而 $Q=\int J^0$ 在复合理论中守恒（由 $\partial_\mu J^\mu=0$ 与 Gauss 定理）。对 $T^1{}_0$、$T^1{}_1$ 的具体化就是 p.110 的两个条件]

例如，世界面能量 $H$ 的守恒要求能动张量的非对角分量
% [原文 p.110]
连续地粘合：
\[
T^1{}_0-\hat T^1{}_0=0
\]
在 $D$ 的各点处成立。满足此条件的缺陷称为共形缺陷 [49]。
% [Tutor Note: p.109–110 处补出 Noether 论证：世界面平移不变性给出 $\partial_\mu T^\mu{}_\nu=0$；对区域 $\Sigma$ 积分并用 Gauss 定理，能量的变化等于通过缺陷 $D$ 的能流 $T^1{}_0$ 的面积分，故 $H$ 守恒当且仅当 $T^1{}_0-\hat T^1{}_0=0$。同理 $T^1{}_1$ 对应沿 $D$ 的动量流；这就是“1-分量连续粘合”的几何含义]

如果缺陷粘合条件还保证能动张量的对角分量连续地粘合：
\[
T^1{}_1-\hat T^1{}_1=0,
\]
那么世界面动量 $P$ 也是守恒荷。满足此条件的缺陷称为拓扑缺陷 [52]，原因在于 $D$ 在世界面上的位置可以光滑地变形而不影响关联函数的值（只要它不越过某个局部算子插入的位置）。

在随后的几节中，我们将写下作用量、变分作用量，然后系统地分析哪些粘合条件极小地置零 $(\delta S)_D$，并保证相关 Noether 流的 1-分量连续地穿过 $D$ 粘合。

\section{玻色 \texorpdfstring{$\sigma$}{σ}-模型的拓扑缺陷}

首先我们分析玻色 $\sigma$-模型的拓扑缺陷。

让我们固定记号。如上，设 $(\sigma^0,\sigma^1)$ 是世界面坐标，在其中缺陷线 $D$ 局部地由 $\sigma^1=0$ 给出，并取世界面度规为平坦的、号差为 $(-,+)$。设 $\Sigma$ 是由 $\sigma^1\ge 0$ 给出的区域，$\hat\Sigma$ 是由 $\sigma^1\le 0$ 给出的区域。带缺陷的玻色 $\sigma$-模型的场由映射 $\Phi:\Sigma\to X$ 与 $\hat\Phi:\hat\Sigma\to\hat X$ 组成；用 $X$ 上的局部坐标 $\phi^i$ 与 $\hat X$ 上的 $\hat\phi^i$，我们可以用函数 $\phi^i(\sigma)$（$\sigma^1\ge0$）与 $\hat\phi^i(\sigma)$（$\sigma^1\le0$）描述这些映射。（注意靶空间被假定具有相同的维数 $n$，因为我们寻找的是拓扑缺陷——而不仅仅是共形缺陷。）同样，乘积映射 $\Phi\times\hat\Phi|_D:D\to X\times\hat X$ 可以由函数 $\phi^I=(\phi^i,\hat\phi^i)|_{(\sigma^0,0)}$ 描述。指标 $i,j$ 取值 $1,\ldots,n$，而 $I,J$ 取值 $1,\ldots,2n$。

带缺陷理论的总作用量是体项之和，再加上一项让体场与活在子流形 $Y\subseteq X\times\hat X$ 上的秩一向量丛之联络耦合的项：
\begin{equation*}
\begin{aligned}
S=&\int_\Sigma \dd^2\sigma\left(-\tfrac12 g_{ij}\partial_\mu\phi^i\partial^\mu\phi^j-\tfrac12 b_{ij}\epsilon^{\mu\nu}\partial_\mu\phi^i\partial_\nu\phi^j\right)\\
&+\int_{\hat\Sigma}\dd^2\sigma\left(-\tfrac12\hat g_{ij}\partial_\mu\hat\phi^i\partial^\mu\hat\phi^j-\tfrac12\hat b_{ij}\epsilon^{\mu\nu}\partial_\mu\hat\phi^i\partial_\nu\hat\phi^j\right)+\int_D \dd\sigma^0 \cA_I\partial_0\phi^I,
\end{aligned}
\end{equation*}
其中 $\cA=\cA_I(\phi)\dd\phi^I$ 是联络 1-形式，$g_{ij}(\phi)$ 与 $b_{ij}(\phi)$ 是 $X$ 上的度规与 B-场，$\hat g_{ij}(\phi)$ 与 $\hat b_{ij}(\phi)$ 是 $\hat X$ 上的度规与 B-场。为简单起见，我们假定 $b$ 与 $\hat b$
% [原文 p.111]
都是闭 2-形式。

变分此作用量并挑出定域在 $D$ 上的项，可得
\begin{equation}
\begin{aligned}
(\delta S)_D=\int_D \dd\sigma^0\Big(&-\left(g_{ij}\partial_1\phi^i+b_{ij}\partial_0\phi^i\right)\delta\phi^j+\left(\hat g_{ij}\partial_1\hat\phi^i+\hat b_{ij}\partial_0\hat\phi^i\right)\delta\hat\phi^j\\
&-\left(\partial_I\cA_J-\partial_J\cA_I\right)\partial_0\phi^I\delta\phi^J\Big).
\end{aligned}
\tag{9.1}
\end{equation}
设 $\sigma=(\sigma^0,0)$ 是 $D$ 上一点，我们来写下在 $\sigma$ 处对场足以保证上述表达式对所有允许变分为零的粘合条件。

正如引言中所提到的，粘合条件的选取对应于子流形 $Y\subseteq X\times\hat X$ 的选取。设 $y=\Phi\times\hat\Phi|_D(\sigma)$ 是 $\sigma$ 在乘积映射下的像。我们首先施加 Dirichlet 条件
\[
y\in Y.
\]
若 $k$ 是 $Y$ 的维数，$0\le k\le 2n$，则上式表示对 $2n$ 个玻色场 $\phi^I$ 的 $2n-k$ 个独立粘合约束。为保证 (9.1) 为零，我们用一组关于导数 $\partial_1\phi^I$ 的 $k$ 个 Neumann 条件来补充这些 Dirichlet 条件。用三个靶空间切向量 $u,v,w\in T_y(X\times\hat X)$ 来描述最为方便，其中
\begin{gather*}
u=u^I\partial_I=(\partial_0\phi^I)\partial_I \\
v=v^I\partial_I=(\partial_1\phi^I)\partial_I \\
w=w^I\partial_I=(\delta\phi^I)\partial_I 
\end{gather*}
这里 $\partial_I\equiv\frac{\partial}{\partial\phi^I}$。约束 $D$ 的点被映到子流形 $Y$ 要求 $u\in TY$，即 $Y$ 的切向子空间；此外，允许变分 $\delta\phi^I$ 是那些使 $w\in TY$ 的变分。
% [Tutor Note: p.111 的计数可核对：$2n$ 个玻色场 $\phi^I$，Dirichlet 条件给出 $2n-k$ 个约束（垂直于 $Y$ 的方向），Neumann 条件给出 $k$ 个约束（沿 $TY$ 的每个独立方向 $w$ 一个），合计 $2n$ 个条件，与场的数目相符；这正是开弦带 B-场时 Dirichlet/Neumann 分解的推广]

这里以及全文中，我们只会有机会讨论在 $Y$ 的点处取值的向量；因此我们把所有向量都视为位于拉回丛 $e^*TM$ 中，其中 $e:Y\to M=X\times\hat X$ 是嵌入映射，并把 $TY$ 视为子丛 $TY\subseteq e^*TM$。

由于靶空间度规 $g$ 与 $\hat g$ 以相反的符号进入 $(\delta S)_D$，我们发现定义 $X\times\hat X$ 上的中性号差度规与 B-场是有用的：
\[
G_{IJ}=\begin{pmatrix}g_{ij}&0\\ 0&-\hat g_{ij}\end{pmatrix},\qquad
B_{IJ}=\begin{pmatrix}b_{ij}&0\\ 0&-\hat b_{ij}\end{pmatrix}.
\]
此外，我们记 $F=-\cF-e^*B$ 为 $Y$ 上的闭 2-形式，它由线丛的曲率 $\cF=\dd\cA$ 与 B-场的拉回组合而成。
